% ---------------------------------------------------------------------------
\begin{frame}{First Normal Form (1NF) --- Definition}
  \begin{block}{Definition}
    A relation is in \textbf{1NF} if and only if:
    \begin{itemize}
      \item Every attribute contains only \textbf{atomic} (indivisible) values
      \item Every entry in a column is of the \textbf{same data type}
      \item There are \textbf{no repeating groups} and no multi-valued attributes
    \end{itemize}
  \end{block}

  \begin{columns}[t]
    \begin{column}{0.48\textwidth}
      \textbf{What Violates 1NF?}
      \begin{itemize}
        \item A \texttt{Songs} column containing
              \texttt{"Song1; Song2"} (two values in one cell)
        \item Arrays or comma-separated lists in a column
        \item Repeating groups: \texttt{Phone1}, \texttt{Phone2}, \texttt{Phone3}
      \end{itemize}
    \end{column}
    \begin{column}{0.48\textwidth}
      \textbf{How to Fix It}
      \begin{itemize}
        \item Give each atomic value its own \textbf{row}
        \item After the fix: one row per (Album, Song) pair
        \item The primary key becomes composite:\\
              (\texttt{AlbumID}, \texttt{Song})
      \end{itemize}
    \end{column}
  \end{columns}

  \vspace{0.4em}
  \begin{alertblock}{Side Effect}
    Applying 1NF makes the table \emph{longer}. Redundancy may \emph{increase}
    temporarily --- this is addressed by 2NF and 3NF.
  \end{alertblock}
\end{frame}

% ---------------------------------------------------------------------------
\begin{frame}{1NF --- Before: Non-Atomic Values}
  \textbf{Before (NOT in 1NF):} multiple song titles crammed into one cell.

  \vspace{6pt}
  \begin{center}
  \small
  \begin{tabular}{lllll}
    \toprule
    \textbf{AlbumID} & \textbf{Title} & \textbf{Songs} & \textbf{Artist} & \textbf{Debut} \\
    \midrule
    A100 & Midnights & Would've\ldots; Karma & Taylor Swift & 2006 \\
    A101 & 1989 (TV) & Shake It Off; Blank Space & Taylor Swift & 2006 \\
    \bottomrule
  \end{tabular}
  \end{center}

  \vspace{6pt}
  \begin{alertblock}{Problems with this design}
    \begin{itemize}
      \item Cannot query for a specific song by name without string parsing
      \item Cannot add a foreign key to \texttt{Songs}
      \item No way to store a featured artist per song
    \end{itemize}
  \end{alertblock}
\end{frame}

% ---------------------------------------------------------------------------
\begin{frame}{1NF --- After: One Song Per Row}
  \textbf{After applying 1NF} --- one row per (Album, Song) pair:

  \vspace{6pt}
  \begin{center}
  \scriptsize
  \begin{tabular}{llllll}
    \toprule
    \textbf{\underline{AlbumID}} & \textbf{\underline{Song}} & \textbf{Featured} & \textbf{Title} & \textbf{Artist} & \textbf{Debut} \\
    \midrule
    A100 & Would've, Could've, Should've & ---       & Midnights & Taylor Swift & 2006 \\
    A100 & Karma                         & Ice Spice & Midnights & Taylor Swift & 2006 \\
    A101 & Shake It Off                  & ---       & 1989 (TV) & Taylor Swift & 2006 \\
    A101 & Blank Space                   & ---       & 1989 (TV) & Taylor Swift & 2006 \\
    \bottomrule
  \end{tabular}
  \end{center}
  \normalsize

  \vspace{6pt}
  \begin{alertblock}{1NF satisfied --- but a new problem appears}
    \texttt{Title}, \texttt{Artist}, and \texttt{Debut} repeat for every song of
    the same album.  The primary key is \texttt{(AlbumID, Song)}, but those
    columns depend only on \texttt{AlbumID} --- a \textbf{partial dependency}.
    This is addressed by \textbf{2NF}.
  \end{alertblock}
\end{frame}

% ---------------------------------------------------------------------------
\begin{frame}{Second Normal Form (2NF) --- Definition}
  \begin{block}{Definition}
    A relation is in \textbf{2NF} if:
    \begin{itemize}
      \item It is in 1NF, AND
      \item Every non-prime attribute is \textbf{fully functionally dependent}
            on the \emph{entire} primary key (no partial dependencies)
    \end{itemize}
  \end{block}

  \begin{block}{Identifying the Partial Dependencies}
    Primary key after 1NF: $\{\texttt{AlbumID}, \texttt{Song}\}$

    \vspace{4pt}
    \small
    \begin{tabular}{lll}
      \toprule
      \textbf{Attribute} & \textbf{Depends on \ldots} & \textbf{Problem?} \\
      \midrule
      \texttt{Featured}  & AlbumID + Song (both)  & No --- fully dependent \\
      \texttt{Title}     & AlbumID only           & \textbf{Partial FD!} \\
      \texttt{Artist}    & AlbumID only           & \textbf{Partial FD!} \\
      \texttt{Debut}     & AlbumID only           & \textbf{Partial FD!} \\
      \bottomrule
    \end{tabular}
  \end{block}

  \begin{alertblock}{Key Insight}
    2NF problems can \emph{only} arise when the primary key is \textbf{composite}.
    A single-attribute PK is automatically in 2NF (if already in 1NF).
  \end{alertblock}
\end{frame}

% ---------------------------------------------------------------------------
\begin{frame}{2NF --- The Fix: Split the Table}
  \textbf{Solution:} One table per independent fact.  Move album-level attributes
  out of the songs table.

  \vspace{0.5em}
  \begin{columns}[t]
    \begin{column}{0.48\textwidth}
      \textbf{Albums} (\underline{AlbumID})
      \small
      \begin{tabular}{llll}
        \toprule
        \textbf{AlbumID} & \textbf{Title} & \textbf{Artist} & \textbf{Debut} \\
        \midrule
        A100 & Midnights & Taylor Swift & 2006 \\
        A101 & 1989 (TV) & Taylor Swift & 2006 \\
        \bottomrule
      \end{tabular}
    \end{column}
    \begin{column}{0.48\textwidth}
      \textbf{Songs} (\underline{AlbumID, Song})
      \small
      \begin{tabular}{lll}
        \toprule
        \textbf{AlbumID} & \textbf{Song} & \textbf{Featured} \\
        \midrule
        A100 & Would've\ldots & --- \\
        A100 & Karma & Ice Spice \\
        A101 & Shake It Off & --- \\
        \bottomrule
      \end{tabular}
    \end{column}
  \end{columns}

  \vspace{0.6em}
  \begin{block}{Status Check}
    \begin{itemize}
      \item \texttt{Songs}: \texttt{Featured} depends on the full composite key \checkmark
      \item \texttt{Albums}: every attribute depends on \texttt{AlbumID} \checkmark
      \item But: \texttt{Debut} still depends on \texttt{Artist}, not \texttt{AlbumID}!
    \end{itemize}
    This \textbf{transitive dependency} is addressed by \textbf{3NF}.
  \end{block}
\end{frame}
